\documentclass{article}

\usepackage{PRIMEarxiv}
\usepackage[utf8]{inputenc}
\usepackage[T1]{fontenc}
\usepackage{cite}
\usepackage{hyperref}
\usepackage{url}
\usepackage{booktabs}
\usepackage{amsfonts}
\usepackage{nicefrac}
\usepackage{microtype}
\usepackage{graphicx}
\usepackage{amsmath}
\graphicspath{{media/}}

\title{A Review of Artificial Intelligence Assisted Antenna Optimization Design}

\author{
  Weiqiang Duan \\
  Course Assignment \\
  Student ID: 12533582
}

\begin{document}
\maketitle

\begin{abstract}
With the rapid development of wireless communication and radar sensing technologies, modern engineering systems impose increasingly demanding requirements on antenna performance. Conventional antenna design methods based on numerical full-wave electromagnetic simulation face several bottlenecks, including high computational cost, long optimization cycles, and a tendency to become trapped in local optima. This paper reviews recent progress in artificial intelligence (AI) assisted antenna design. First, the applications of artificial neural networks (ANNs), deep neural networks (DNNs), and graph neural networks (GNNs) in antenna performance prediction, complex-structure representation, and array synthesis are summarized. Second, the development of cooperative optimization frameworks that combine meta-heuristic algorithms with deep-learning surrogate models is discussed. The reviewed literature shows that AI methods can significantly improve nonlinear modeling capability and multi-objective optimization efficiency in complex antenna design. However, important limitations remain, including dependence on high-quality data, weak cross-structure generalization, limited physical interpretability, and insufficient verification of manufacturability. This review further identifies research gaps in standardized benchmark datasets, deep integration of electromagnetic physics with data-driven learning, and system-level multiphysics co-optimization. Finally, future AI antenna design is argued to move beyond fragmented one-task-one-model pipelines toward antenna foundation models trained on large-scale electromagnetic data, full-domain electromagnetic topology learning with graph deep learning, and closed-loop intelligent design platforms integrating requirement analysis, structure generation, simulation, and feedback.
\end{abstract}

\keywords{AI-assisted design \and deep-learning surrogate model \and graph neural network \and antenna foundation model \and multi-objective optimization}

\section{Introduction}

The rapid growth of wireless communication, radar sensing, the Internet of Things, and fifth- and sixth-generation mobile systems has placed increasingly strict requirements on antenna performance. Multiband operation, ultrawide bandwidth, high gain, low sidelobe level, and compact size have become common design objectives. At present, full-wave electromagnetic simulation tools such as HFSS and CST are widely used in antenna analysis and optimization. These tools are indispensable for accurate electromagnetic evaluation, but they also introduce significant practical limitations. Full-wave solvers require substantial computational resources, and trial-and-error parameter tuning depends heavily on expert experience. When many structural parameters must be swept or optimized, the process is time-consuming and may easily converge to local optima rather than globally satisfactory designs \cite{sarker2023review,misilmani2020review,khan2022review}.

To reduce computational cost and shorten the design cycle, deep-learning methods represented by artificial neural networks (ANNs) and deep neural networks (DNNs) have gradually been introduced into antenna design. As early as 1998, Mishra and Patnaik demonstrated the feasibility of neural-network-based computer-aided design for square patch antennas \cite{mishra1998square}. Since then, deep learning has become attractive because it can approximate highly nonlinear relationships between antenna structures and electromagnetic responses. By constructing lightweight surrogate models, these methods can quickly predict the response of unseen antenna structures and greatly reduce the time required for performance evaluation \cite{li2019deepnis,yao2018fdtd,zhang2003ann,lecun2015deep,shi2018emc,feng2022ann}.

However, forward prediction alone is not sufficient for inverse antenna design. In a large, non-convex design space, a key problem is how to infer structural parameters that satisfy target specifications. Meta-heuristic algorithms, including genetic algorithms (GA), particle swarm optimization (PSO), differential evolution, and newer swarm-intelligence methods, therefore play an important role. These algorithms do not require explicit gradient information and are suitable for global search under complex constraints and multiple objectives. Recent studies increasingly adopt a cooperative framework in which a surrogate model performs fast performance evaluation while a meta-heuristic optimizer updates candidate structures. This "surrogate evaluation plus global search" strategy balances computational efficiency and global optimization capability and has become an important route in AI-assisted antenna design \cite{singh2024polygon,huang2024meta,pietrenko2023resolution,yang2026gphercnn}.

Although this cooperative framework works well for single antennas and regular arrays, it remains limited when applied to modern antenna systems with nonuniform sparse arrays, conformal arrays, and diverse topologies. Traditional ANNs and DNNs generally operate on fixed-dimensional Euclidean tensors, which makes it difficult to describe irregular spatial layouts and complicated electromagnetic coupling. Graph neural networks (GNNs) provide a promising solution because they model non-Euclidean graph-structured data. In antenna array problems, array elements, physical distances, feeding relations, and coupling interactions can be represented as graph nodes and edges. Therefore, GNNs can break the limitation of fixed topology and provide high-fidelity modeling of complex array structures \cite{yang2024doa,wu2021gnn,yu2018stgcn,he2025fluid,xie2025pinching}.

Modern AI-assisted antenna design is therefore not the isolated application of a single algorithm. Instead, it is a design paradigm jointly driven by deep-learning feature extraction, meta-heuristic global optimization, and graph-based topology modeling. A central future challenge is to integrate geometric structure, electromagnetic response, multi-objective constraints, and intelligent search into a more general and closed-loop framework, thereby improving generalization for complex antenna systems.

From the perspective of an engineering workflow, this shift also changes the role of electromagnetic simulation. In conventional design, full-wave simulation is the main evaluator in almost every iteration. In an AI-assisted workflow, simulation gradually becomes a source of high-quality data, a verifier for promising candidates, and a feedback mechanism for model correction. This does not mean that AI replaces electromagnetic theory. Rather, AI can be viewed as an acceleration layer built on top of physical modeling, numerical simulation, and experimental validation. A useful review should therefore examine not only whether a model achieves high prediction accuracy, but also how it is embedded into the entire antenna design process.

Accordingly, this paper organizes the literature around three technical questions. The first is how deep-learning surrogate models predict antenna responses efficiently. The second is how meta-heuristic algorithms use these fast predictors to search for better structures. The third is how graph-based models describe complex topologies that cannot be easily represented by fixed-length parameter vectors. These three questions correspond to prediction, optimization, and representation, which together form the core of AI-assisted antenna design.

\section{Literature Review}

\subsection{Deep Learning for Antenna Design}

Traditional antenna design and optimization rely heavily on electromagnetic theory, engineering experience, and repeated parameter sweeps in simulation software. This process is often slow and does not guarantee global optimality. Deep learning offers a new route by learning the mapping between antenna structural parameters and performance metrics. Once trained, the model can replace part of the repeated full-wave simulation workload and support rapid prediction and optimization.

ANNs and DNNs are among the most mature AI models in antenna design. They have been used for structure design, performance prediction, and array synthesis. The early work of Mishra and Patnaik showed that neural networks can be used to build a CAD model for square patch antennas \cite{mishra1998square}. Later studies improved model structures and training strategies to better fit antenna design tasks. Reviews of antenna-oriented machine learning also show that neural networks are effective tools for mapping geometry, feeding conditions, and electromagnetic responses \cite{misilmani2020review,feng2022ann}.

Recent studies focus more on low-sample surrogate modeling, automatic structure generation, and multi-objective inverse design. Belen et al. proposed a low-cost surrogate modeling method for transmitarray antennas and used automated deep learning to design transmitarray units in the 8--14 GHz, 22--28 GHz, and 28--36 GHz bands with a small training dataset \cite{belen2023transmitarray}. Wei et al. developed a CNN-LSTM based automatic modeling framework that can generate antenna modeling code from antenna images, thereby accelerating sample generation for different antenna structures \cite{wei2023cnn}. Singh et al. combined a CNN surrogate model with multi-island differential evolution for multi-objective inverse design of polygon patch antennas, optimizing reflection coefficient, input impedance, and radiation pattern within a short time \cite{singh2024polygon}.

Deep learning has also been used for antenna array design and synthesis. Ayestaran et al. used neural networks to synthesize nonuniform antenna arrays by relating radiation patterns to element excitation voltages, enabling array synthesis without strong assumptions about geometry \cite{ayestaran2007nonuniform}. Mishra et al. used radial basis function (RBF) neural networks to estimate first null beamwidth for broadside and end-fire uniform linear arrays, showing that RBF networks can adapt to nonlinear array behavior \cite{mishra2015rbf}.

Overall, deep-learning methods offer strong nonlinear modeling capability and can improve the efficiency of antenna structure design, performance prediction, and array synthesis. Nevertheless, conventional ANN and DNN models are mainly suitable for Euclidean data and fixed topologies. Their limitations in irregular arrays and non-Euclidean structures motivate the use of graph deep learning.

Another important issue in deep-learning-based antenna design is the choice of input and output representation. Some studies use explicit geometric parameters such as patch length, substrate height, feed position, or slot dimensions as inputs. This representation is compact and easy to train, but it is usually tied to one antenna family. Other studies use images, meshes, or field samples as inputs, which allows the model to capture richer structural information but requires larger datasets and more careful preprocessing. Similarly, output labels may include S-parameters, resonant frequency, gain, radiation efficiency, or full far-field patterns. The more complete the output representation is, the more useful the model becomes for system-level design, but the learning problem also becomes more difficult.

The training strategy also affects practical performance. If samples are generated by uniform parameter sweeps, many simulations may be wasted on unimportant regions of the design space. Active learning and adaptive sampling can reduce this waste by selecting samples near performance boundaries or uncertainty regions. Low-fidelity simulations may also be used in early training, while high-fidelity simulations refine the model near the final optimum. These strategies show that data generation is not merely a preparation step; it is part of the design methodology itself.

\subsection{Graph Deep Learning for Antenna Design}

Graph neural networks are designed to process data represented by nodes and edges. They are useful when the relationship among objects is as important as the attributes of the objects themselves. GNNs have been widely studied in recommendation systems, traffic forecasting, biological modeling, and circuit analysis. The spatio-temporal graph convolutional network proposed by Yu et al., for example, models traffic networks by combining graph structure and temporal dynamics \cite{yu2018stgcn}. Such work provides methodological inspiration for electromagnetic and antenna problems.

In antenna systems, array elements, users, feed points, and movable antenna positions naturally form relational structures. Therefore, GNNs can describe element connectivity, spacing, and coupling effects. Yang et al. proposed a residual graph convolutional network based on GraphSAGE for direction-of-arrival estimation in random sparse linear arrays. By constructing received signals as graph data and aggregating neighbor information, the method compensates for missing element information and improves robustness under low signal-to-noise ratio and low snapshot conditions \cite{yang2024doa}.

GNNs are also emerging in reconfigurable antenna and communication systems. He et al. proposed a two-stage GNN framework for fluid antenna systems, separately inferring antenna positions and beamforming vectors and training the model with an unsupervised loss function to optimize sum rate and energy efficiency \cite{he2025fluid}. Xie et al. modeled a pinching-antenna downlink system as a bipartite graph and proposed a graph-attention-based model for joint antenna placement and power allocation \cite{xie2025pinching}.

These studies show that GNNs offer a structural representation that is more suitable for complex antenna arrays than traditional fixed-input networks. By incorporating element positions, distances, coupling relations, and topology into one graph model, GNNs can overcome the limitations of regular-array assumptions and support more flexible antenna synthesis and optimization.

For antenna arrays, the meaning of graph nodes and edges can be defined in different ways. A node may represent an antenna element, a feeding port, a user, or a movable radiation point. An edge may represent physical distance, electromagnetic coupling strength, signal correlation, or a communication link. This flexibility is valuable because different antenna systems emphasize different relations. For example, a sparse array used for direction finding may focus on missing elements and spatial sampling, while a fluid antenna system may focus on candidate positions and beamforming variables. A GNN framework can incorporate these relations without forcing all systems into the same rectangular tensor format.

At the same time, GNN-based antenna design is still at an early stage. Many current studies test graph models on specific array scenarios, and the transferability among different array geometries remains unclear. The construction of graph edges is also a nontrivial problem. If edges are defined only by distance, important coupling mechanisms may be missed; if all possible couplings are included, the graph may become dense and computationally expensive. Future work needs to combine electromagnetic insight with graph construction so that the topology is physically meaningful rather than merely convenient for machine learning.

\subsection{Meta-Heuristic Algorithms and Surrogate-Assisted Optimization}

In AI-assisted antenna design, machine-learning and deep-learning models are often used as surrogate models for fast prediction, while meta-heuristic algorithms are used for inverse optimization. Meta-heuristic algorithms simulate natural or collective search behavior and are well suited for nonlinear, multivariable, and multi-objective problems. GA and PSO are classical examples, and newer algorithms such as seagull optimization, butterfly optimization, grey wolf optimization, and coati optimization have also been explored.

Recent research increasingly combines surrogate models with intelligent optimizers. In the polygon patch antenna case study, a CNN surrogate model was used as a fast evaluator, while multi-island differential evolution performed multi-objective search \cite{singh2024polygon}. Huang et al. investigated antenna modeling based on meta-heuristic optimization of neural-network weights and biases, comparing several swarm-intelligence methods and showing that improved butterfly optimization can enhance prediction accuracy and efficiency \cite{huang2024meta}.

To address the high cost of electromagnetic simulation, Pietrenko-Dabrowska et al. studied expedited meta-heuristic antenna optimization using electromagnetic model resolution management, switching among models of different fidelities to reduce computational burden \cite{pietrenko2023resolution}. Yang et al. further combined a 1D GP-HetCNN surrogate with WOA, COA, GWO, PSO, and Bayesian optimization for ultrawideband antenna parameter optimization, using only a small number of HFSS simulations for verification \cite{yang2026gphercnn}. These works indicate that surrogate models, multifidelity modeling, and meta-heuristic algorithms are becoming increasingly integrated.

The main strength of meta-heuristic algorithms is that they can search complex spaces without requiring differentiable objective functions. This is especially useful in antenna design because design objectives often include discontinuous constraints, fabrication rules, and multiple performance indicators. However, this flexibility comes with a cost: a population-based optimizer may require thousands of evaluations. If every evaluation calls a full-wave solver, the total cost becomes unacceptable. Surrogate-assisted optimization reduces this burden by allowing the optimizer to evaluate many candidates quickly and reserve expensive simulations for verification or model updating.

In practice, the performance of a surrogate-assisted optimizer depends on the balance between exploration and reliability. If the optimizer trusts an inaccurate surrogate too much, it may converge to a false optimum that performs well in prediction but poorly in simulation. If the process relies too frequently on high-fidelity simulation, the efficiency advantage is weakened. A robust workflow should therefore include error monitoring, periodic electromagnetic validation, and possibly uncertainty estimation. Such mechanisms can prevent the optimizer from exploiting artifacts of the surrogate model.

\section{Discussion}

\subsection{Advantages of Existing AI-Based Methods}

The most direct advantage of AI-assisted antenna design is improved efficiency. Conventional design requires repeated full-wave simulations, whereas a trained surrogate model can rapidly predict antenna performance and reduce the number of expensive simulations \cite{sarker2023review,khan2022review}. In low-cost transmitarray design, for example, a relatively small dataset was sufficient to construct a useful surrogate model and validate multiple frequency bands experimentally \cite{belen2023transmitarray}.

AI methods also improve nonlinear modeling and multi-objective optimization. Antenna performance is highly sensitive to geometry, feeding position, element spacing, and material properties. Neural networks, RBF models, CNNs, and hybrid architectures can approximate such nonlinear relationships more flexibly than simple analytical formulas \cite{feng2022ann,mishra2015rbf}. When combined with meta-heuristic algorithms, surrogate models can support simultaneous optimization of reflection coefficient, input impedance, radiation pattern, and bandwidth \cite{singh2024polygon,huang2024meta}.

A third advantage is the representation of complex structures. CNN-LSTM models can process antenna images and generate modeling code \cite{wei2023cnn}. GNNs can directly represent array elements and their relations as graph nodes and edges, which is particularly valuable for sparse arrays, fluid antenna systems, and pinching-antenna systems \cite{yang2024doa,he2025fluid,xie2025pinching}.

These advantages also interact with each other. Efficient prediction makes it possible to test more candidate structures; better nonlinear modeling makes the prediction reliable enough to guide optimization; and richer structural representation expands the range of antenna systems that can be optimized. For example, a design framework for a nonuniform array may use a GNN to encode topology, a neural surrogate to estimate radiation performance, and a meta-heuristic algorithm to adjust excitation or element locations. The value of AI-assisted design is therefore not only in a single accurate model, but in the ability to connect multiple computational modules into a faster design loop.

Another practical advantage is knowledge reuse. Once a model has learned useful representations from one class of antennas, part of that knowledge may be transferred to related structures. Although current transfer learning in antenna design is still limited, the idea is important because many engineering projects modify an existing antenna rather than design a completely new one. If a model can reuse information from earlier simulations, the cost of adapting to a new frequency band or a slightly different geometry can be reduced.

\subsection{Limitations}

Despite these advantages, AI-based antenna design still depends strongly on training data. High-quality antenna datasets are usually generated by full-wave simulation, which means that part of the computational burden is shifted from the optimization stage to the dataset construction stage \cite{sarker2023review,belen2023transmitarray}. For high-dimensional structures such as large arrays and metasurfaces, data requirements can become a major bottleneck.

Generalization and physical interpretability are also unresolved. Many ANN and CNN surrogate models learn statistical patterns from a specific dataset and may perform poorly when antenna topology, frequency band, material, or boundary conditions change. Moreover, black-box models do not guarantee consistency with Maxwell's equations, energy conservation, or reciprocity. Physics-informed neural networks and physics-informed machine learning provide important inspiration by embedding differential equations or physical constraints into the learning process \cite{raissi2019pinn,karniadakis2021piml}.

Engineering verification and manufacturability remain additional challenges. AI optimization can produce ideal numerical structures that are difficult to fabricate or sensitive to tolerance, dielectric loss, assembly error, or feed-network complexity. Although some studies include experimental validation \cite{belen2023transmitarray}, many remain at the simulation or algorithm-comparison stage. A practical AI design workflow should include manufacturing constraints, anechoic-chamber testing, and feedback from measured results.

A related limitation is that many studies evaluate models only by numerical accuracy, such as mean squared error of S-parameters or prediction error of resonant frequency. These metrics are useful, but they do not fully describe engineering value. A model with slightly higher prediction error may still be more useful if it preserves the correct trend of parameter changes or consistently identifies high-quality candidate regions. Conversely, a model with low average error may fail near sharp resonances or constraint boundaries, where design decisions are most sensitive. Evaluation should therefore include both statistical metrics and design-oriented metrics.

There is also a gap between academic datasets and industrial antenna development. In real projects, data may be confidential, incomplete, noisy, or generated under different simulation settings. Models trained on clean synthetic data may not handle these conditions well. In addition, industrial workflows often require traceability: engineers need to know why a design was selected, what constraints were active, and how robust the final structure is. These requirements increase the importance of interpretable models, uncertainty estimation, and standardized reporting.

\subsection{Research Gaps}

First, the field lacks standardized datasets and evaluation protocols. Different studies use different antenna types, frequency bands, simulation tools, mesh settings, and metrics. This makes fair comparison among neural networks, GNNs, and meta-heuristic optimizers difficult \cite{sarker2023review,yang2026gphercnn}. Open benchmark datasets covering multiple bands and topologies would support reproducible research.

Second, the integration of physics and data-driven learning is still insufficient. Most models remain pure function approximators, even though antenna behavior is governed by electromagnetic field distributions and boundary conditions. Embedding electromagnetic priors, equivalent-circuit knowledge, or partial differential equation constraints into neural architectures could improve small-sample accuracy and cross-band generalization \cite{raissi2019pinn,karniadakis2021piml}.

Third, system-level and multiphysics co-optimization is underdeveloped. Real antenna systems must satisfy electromagnetic, thermal, mechanical, environmental, and communication-system constraints. Han et al. showed that sequential multiphysics machine-learning-assisted optimization can support low-wind-load broadband dual-polarized antenna and array design \cite{han2025multiphysics}. Extending AI antenna design from isolated electromagnetic optimization to system-level multiphysics optimization is therefore an important research direction.

In addition, there is a lack of common workflow standards for closed-loop AI antenna design. Many papers describe a model or optimizer, but fewer specify how data are selected, how simulation fidelity is controlled, how errors are corrected, and how final designs are verified. Without such workflow descriptions, it is difficult to reproduce results or compare methods fairly. A mature research field should provide not only models but also repeatable pipelines, including dataset construction, training protocol, optimizer settings, validation strategy, and reporting format.

\section{Future Directions}

\subsection{From Task-Specific Surrogates to Antenna Foundation Models}

Most current surrogate models are trained for a specific antenna structure, frequency range, or performance metric. This one-task-one-model paradigm performs well within the training distribution but requires new sampling and retraining when the structure or operating condition changes. A promising direction is to build antenna foundation models trained on large-scale, multi-type, and multimodal electromagnetic data. Such models could learn shared representations among geometry, material parameters, feeding methods, field distributions, S-parameters, and far-field patterns. With limited fine-tuning data, they could then adapt to new antenna structures or design tasks, improving transferability and reducing sampling cost \cite{tang2024survey}.

An antenna foundation model would require data far beyond simple tabular parameter-response pairs. Useful pretraining data may include CAD geometry, mesh information, material distribution, near-field maps, far-field patterns, S-parameter curves, simulation boundary settings, and even textual design requirements. Learning from these heterogeneous sources could allow the model to understand both structure and function. For example, the model may learn that certain slot shapes tend to introduce resonances, or that element spacing strongly affects sidelobe behavior. Such knowledge could then be reused when a new design problem is introduced.

However, foundation models also raise practical challenges. Large-scale electromagnetic datasets are expensive to generate, and different simulation tools may produce data in different formats. Model size, training cost, and verification difficulty must also be considered. Therefore, near-term progress may come from domain-specific pretraining on moderate-scale datasets, followed by transfer learning, retrieval-based design assistance, and active learning. Even such intermediate steps would represent a significant improvement over isolated task-specific surrogates.

\subsection{Extending Graph Learning to Full-Domain Electromagnetic Topologies}

Graph deep learning can be extended beyond sparse arrays to circular arrays, conformal arrays, metasurfaces, metamaterial antennas, and other electromagnetic structures. Future work should construct more general graph representations for electromagnetic problems, including nodes for physical elements, edges for geometric or coupling relations, and attributes for material or field information. This could allow GNNs to serve not only as performance predictors but also as independent tools for array synthesis and topology optimization \cite{su2024gnn}.

A full-domain electromagnetic topology model should also support multiple levels of abstraction. At a low level, a graph may describe mesh cells, current paths, or local material regions. At a middle level, it may describe antenna elements, feeding networks, and coupling paths. At a high level, it may describe system components such as arrays, users, reflectors, or reconfigurable surfaces. Connecting these levels would make it possible to reason about both local electromagnetic behavior and system-level performance. This is especially relevant for future intelligent surfaces and integrated sensing and communication systems, where antenna design cannot be separated from the surrounding propagation environment.

\subsection{Toward Closed-Loop Intelligent Design Platforms}

Current AI antenna workflows still consist of loosely connected modules, including structure modeling, data generation, surrogate training, intelligent optimization, full-wave simulation, and result correction. A more advanced platform should connect these stages into a closed loop. Given requirements such as operating frequency, gain, bandwidth, polarization, size, and application scenario, the system could automatically generate candidate structures, predict performance, call an optimizer, run HFSS or CST simulations, analyze results, and update the model. Large language model agents may help schedule tools, generate simulation scripts, interpret results, and select valuable new samples for active learning. This "requirement--generation--simulation--feedback" loop would improve automation, verification, and engineering applicability \cite{fan2024surrogate}.

Such a platform should not be treated as a fully automatic black box. Instead, it should support human-in-the-loop decision making. Engineers may specify hard constraints, reject impractical structures, adjust optimization priorities, or inspect why a candidate was recommended. The platform should record the design history so that each iteration can be traced. This is important for engineering responsibility, especially when an AI-generated antenna is used in communication, radar, aerospace, or other safety-sensitive applications.

Another promising component is active learning. After each simulation or measurement, the platform can identify where the model is uncertain and select the next most informative candidate. This can reduce unnecessary simulations while improving model accuracy in critical regions of the design space. Combined with uncertainty-aware surrogate models, active learning can form a feedback loop in which the model becomes more reliable as the design progresses.

\section{Conclusion}

This paper reviewed the evolution of AI-assisted antenna design, focusing on deep-learning surrogate models, graph neural networks, and meta-heuristic optimization. The literature indicates that AI methods can significantly improve design efficiency, nonlinear modeling ability, complex-structure representation, and multi-objective optimization. At the same time, important challenges remain: high data requirements, weak generalization, limited physical interpretability, and insufficient manufacturability verification.

Future AI antenna design should move beyond fragmented task-specific models toward general, transferable, and physically informed design systems. Antenna foundation models, graph-based full-domain electromagnetic topology learning, physics-informed constraints, active learning, and closed-loop intelligent design platforms are likely to be central to this transition. These developments may help transform AI antenna design from a simulation-assisted technique into a practical end-to-end engineering workflow.

In summary, the future value of AI in antenna design will depend on whether it can connect computational efficiency with physical trustworthiness. Fast prediction alone is not enough, and automatic optimization alone is not enough. The most useful systems will be those that combine accurate electromagnetic knowledge, robust data-driven learning, efficient global search, and verifiable engineering feedback. When these components are integrated carefully, AI-assisted antenna design can become not only faster, but also more systematic, reusable, and reliable.

\bibliographystyle{unsrt}
\bibliography{references}

\end{document}
